Im ersten Schritt (isotherm) gilt $p_1V_1 = p_2V_2$ mit $V_2 = \frac{1}{2}V_1$:
\begin{align}
 p_2 &= p_1\cdot\frac{V_1}{V_2} = 2\,p_1
\end{align}
Im zweiten Schritt (isochor) gilt $p_2/T_2 = p_3/T_3$ mit $T_3 = 2\,T_2$ (absolute Temperatur):
\begin{align}
 p_3 &= p_2\cdot\frac{T_3}{T_2} = 2\,p_2 = \result{4\,p_1}
\end{align}
Der Druck wird viermal so gross. Die Stoffmenge (ein Mol) spielt keine Rolle, solange sie gleich bleibt.
